\chapter*{Preface}
\addcontentsline{toc}{chapter}{Preface}

Understanding the Fourier transform changes how you see a signal. Before, a sound is just a
wiggly line on a screen, an unintelligible scribble of voltage against time. Afterwards
you can \emph{see} that the scribble is really a chord -- a stack of pure tones,
each with its own loudness and its own timing -- and you can pick the chord apart note
by note, push some notes down, lift others up, and put the sound back together changed.
The wiggly line itself has not changed; you have learned to look at
the same object from a second direction, from which its structure is
obvious.

This book is about that second direction. It is about \emph{frequency}: what it really
is, why it is hiding inside every signal, how to measure it, how to manipulate it, and
how to build hardware that does all of this in real time on a battery. The promise of
the title -- \emph{The Geometry of Frequency} -- is meant literally. We are going to
treat signals as \emph{vectors}, frequencies as \emph{directions}, the Fourier transform
as a \emph{rotation of your coordinate system}, and filtering as \emph{reshaping a shadow}.
Almost everything that is usually presented as a wall of integrals and summations can
be drawn as a picture. We will draw the picture first, every time, and only then write down
the formula that makes the picture precise.

\chapterline

\section*{What this book is, and what it is not}

This is a book \emph{from the ground up}: it starts from first principles but does not
stay shallow. We start from the ground -- from a voltage
wiggling on a wire, from a number arriving once every millisecond from a sensor -- and we
build, step by step, all the way up to spectrograms, digital filters, the FFT,
sampling theory, and the signal-processing pipeline inside a real wearable device. By the
end you should be able to read a research paper on gesture recognition, design a bandpass
filter that actually does what you need, and write the C that runs it on a microcontroller.

It is \emph{not} a formula reference. There are excellent ones already, and we will point
to them. If all you want is the equation for the discrete Fourier transform, you can find
it in thirty seconds; you do not need eight hundred pages for that. What eight hundred
pages can buy you is \emph{understanding} -- the kind that lets you re-derive the equation
when you have forgotten it, debug it when it misbehaves, and recognize the same geometric idea
in a new setting three chapters later. Memorized formulas fade; a geometric picture lasts. If you
understand that the Fourier transform projects a signal onto a family of rotating arrows,
then the formula is just bookkeeping, and you will never be truly lost.

\begin{focusbox}
The single most important sentence in this book: \textbf{a signal is a vector, and
frequency analysis is choosing a better basis to look at it in.} Almost every chapter is
a variation on that theme. If a derivation ever feels like meaningless symbol-pushing,
stop and ask: \emph{what vector are we looking at, and what basis are we measuring it in?}
The answer usually makes the derivation clear.
\end{focusbox}

\section*{Who you are}

I am assuming you are smart, curious, and impatient with hand-waving -- and also that you
do \emph{not} already know signal processing. You should be comfortable with the idea of a
function, with basic trigonometry (sine and cosine), and with the notion that you can add
and scale things. Calculus and linear algebra will appear, but I will reintroduce every
piece I lean on, in the language of pictures, at the moment I need it. If you have seen an
integral before, good; if it has gone hazy, the \textbf{Connection} boxes will pull you back
to solid ground. You do not need to have touched a soldering iron, but I hope that by the
end you will want to.

A word on prerequisites in the other direction: do not let \emph{too much} prior exposure
get in your way either. If you were taught Fourier analysis as a sequence of theorems to
survive, you may have to unlearn the reflex of reaching for a formula before you have a
picture. Be patient with the slow, visual openings of each chapter; the later formulas depend
on them.

\section*{The through-line: a wrist gesture controller}

A book that is only a tour of techniques is exhausting to read and easy to forget. So this
book has a destination -- a single concrete machine that we build piece by piece from start
to finish. The machine is a \textbf{wrist- and forearm-worn gesture controller for playing
instruments}. You wear it like a watch or a cuff; it watches how your hand and arm move and
tense; and it emits \textbf{MIDI} -- the universal language of electronic instruments -- so
that a flick of the wrist bends a note, a tap triggers a drum, a slow rotation sweeps a
filter on a synthesizer. The dream is to play music with gesture, with no keyboard between
you and the sound.

Why this goal? Because it touches almost everything in the book. To build it you must
understand \emph{sampling} (how fast to read the sensors), \emph{filtering} (separating the
gesture from the tremor and the noise), the \emph{Fourier and short-time Fourier transforms}
(what frequencies live in a motion, and how they change over time), \emph{convolution and
LTI systems} (the mathematics of every smoothing and shaping operation you will perform),
\emph{feature extraction and template matching} (turning a stream of numbers into the
decision ``that was a tap'' by comparing it against patterns you recorded during a short
calibration), and finally \emph{embedded DSP} (making it all run in a few kilobytes on a
chip that costs a couple of dollars). Each chapter, where it is natural,
ends with a \textbf{Toward the wrist tracker} box that pulls the chapter's idea toward this
build, so that each technique is tied to a concrete use.

\begin{goalbox}
Here is the machine we are building, named once so you can hold it in your head for eight
hundred pages. A small, low-power board on the wrist and forearm carries an \textbf{IMU} (a
high-rate accelerometer plus gyroscope), an array of \textbf{force-myography} (FMG) pads or
force-sensitive resistors that feel the forearm tendons bulge as the fingers move, and a
second high-rate accelerometer pressed to the skin to pick up the faint bio-acoustic taps of
contact. Some versions add a surface-\textbf{EMG} channel for the muscles' electrical pulses.
The wrist node stays deliberately thin: it samples, lightly filters, compresses, and streams
the data over Bluetooth (BLE) to a \textbf{phone}, which handles the heavier computation -- running a
short guided calibration, building a per-user template for each finger movement, and
conditioning the streams. A lightweight template match runs on the device itself for the
low-latency live-play path, emitting MIDI to whatever instrument you like. By the capstone
chapter you will have seen every block of this pipeline derived from first principles and
built in code.
\end{goalbox}

Figure~\ref{ch00:arch} sketches that split. The wrist node does as little as it can get away
with -- read, filter lightly, compress, send -- so it stays small and runs for a long time on
a coin cell. The phone, which has power and memory to spare, runs the calibration and the
deeper analysis, while a small template matcher on the wrist keeps the live-play latency low.

\begin{figure}[H]
\centering
\begin{tikzpicture}[
  block/.style={draw, rounded corners, align=center, font=\scriptsize,
                minimum height=8mm, minimum width=18mm, fill=blue!4},
  sens/.style={draw, rounded corners, align=center, font=\scriptsize,
               minimum height=6mm, minimum width=20mm, fill=green!5},
  every node/.style={inner sep=2pt},
  arr/.style={-{Latex[length=2mm]}, thick},
]
\node[sens] (imu) at (0, 0.8) {IMU\\(accel + gyro)};
\node[sens] (fmg) at (0, 0.0) {FMG / FSR array};
\node[sens] (mic) at (0, -0.8) {contact accel\\(bio-acoustic)};

\node[block] (front) at (2.4, 0.0) {sample +\\light filter};
\node[block] (match) at (4.8, 0.0) {on-device\\template match};
\node[block] (midi)  at (7.2, 0.0) {MIDI out\\(notes / CC)};

\node[block] (ble)   at (4.8, 1.5) {BLE\\stream};
\node[block, fill=purple!4] (phone) at (10.0, 1.5) {phone:\\calibration +\\templates};

\draw[arr] (imu.east) -- (front.west);
\draw[arr] (fmg.east) -- (front.west);
\draw[arr] (mic.east) -- (front.west);

\draw[arr] (front.east) -- (match.west);
\draw[arr] (match.east) -- (midi.west);

\draw[arr] (front.north) |- (ble.west);
\draw[arr] (ble.south) -- node[font=\tiny, right]{templates} (match.north);

\draw[arr] ([yshift=1.5mm]ble.east) -- node[font=\tiny, above]{data stream} ([yshift=1.5mm]phone.west);
\draw[arr] ([yshift=-1.5mm]phone.west) -- node[font=\tiny, below]{templates} ([yshift=-1.5mm]ble.east);

\node[draw, dashed, gray, rounded corners, fit=(imu)(fmg)(mic)(front)(ble)(match)(midi),
      inner sep=4mm, label={[font=\scriptsize\itshape, gray]above:thin wrist node}] {};
\node[draw, dashed, purple!60!black, rounded corners, fit=(phone),
      inner sep=3mm, label={[font=\scriptsize\itshape, purple!60!black]above:companion phone}] {};
\end{tikzpicture}
\caption{The capstone device, split across two computers. The wrist node senses three
channels, conditions them lightly, and streams them over Bluetooth; the phone runs the
per-finger calibration and the heavy offline analysis; a small template matcher on the wrist
gives the live-play path its low latency and drives the MIDI output. No machine learning is
involved -- detection is a direct match against patterns you record during calibration.}
\label{ch00:arch}
\end{figure}

\section*{A straight contract about sensors}

Here is where I want to be completely straight with you, because the field is full of
breathless marketing and it will save you months of frustration to know what is and is not
possible early.

\textbf{Which finger you moved is detectable from the wrist and forearm.} When you flex a
finger, the tendons that drive it slide and the muscles in your forearm bulge in a pattern
that is specific to \emph{that} finger. A band of force-myography (FMG) pads pressed against
the forearm feels those bulges directly; a surface-EMG electrode picks up the electrical
pulses that command the same muscles. Both signals carry a usable per-finger signature.
WristFlex demonstrated reliable detection of which finger pinched from a forearm FMG band;
the Myo armband and the open MyoWare boards put sEMG finger sensing into many hands. So
finger \emph{movement} detection is a solved-enough problem to build
on, provided you do one thing.

That one thing is \textbf{per-finger calibration}. The tendon and muscle signatures differ
from person to person and shift if the band rotates, so we do not ship a fixed model. We run
a short guided enrollment: the app asks you to perform each finger movement a few times, the
device records the FMG/EMG signature each time, and the phone builds a \emph{template} (and
a threshold) for each finger from \emph{your} signals. Live detection is then a direct
comparison -- correlate the incoming signal against the stored templates and fire the
nearest match. No machine learning, no opaque trained model: just your own recorded patterns
and a distance measure. Calibration is therefore the central step in using the instrument.

\begin{trapbox}
The most expensive mistake in wearable sensing is designing the \emph{software} for
capabilities the \emph{hardware} cannot provide -- but the opposite error is just as costly:
under-reaching, and assuming the wrist can only see ``the whole hand moved.'' It can see more
than that, if you sense the right channel (tendons via FMG, muscle pulses via EMG) and
calibrate per user. Throughout this book, whenever a technique brushes against the limits of
what a sensor can physically measure, I will say so plainly, and point to the real published
work so you can check me.
\end{trapbox}

One limit, however, is set by physics.
\textbf{Continuous, per-joint finger angles cannot be recovered from a wrist IMU.} The bulk
motion of the wrist simply does not contain the angle of each knuckle; the problem is
underdetermined, and no clever software conjures information that never entered the device.
If you need the actual joint angles -- a full hand pose, every degree continuously -- you
need sensors on the fingers themselves: instrumented gloves or rings with their own IMUs, as
in the finger-worn ``Learning on the Rings'' line of work. That is a clean upgrade path, and
we will sketch it. But everything our instrument needs -- knowing which finger moved, how
hard, and when it tapped -- lives comfortably on the doable side of that line.

Our build improves on the single-modality lab papers through
\emph{multimodal fusion} rather than a smarter algorithm. A high-rate IMU, an FMG/FSR array, and a second accelerometer used
as a contact and bio-acoustic sensor (the trick behind ViBand) each see a different slice of
the same gesture, and together they disambiguate what any one alone would miss -- the way
PLOS One (2020) combined IMU with acoustics, and MDPI Sensors (2023) counted gestures from a
wrist IMU. Fuse those streams, calibrate per finger, and you have something more capable than
the HR/steps/sleep wearables on the market today, which never attempt finger input at all.
The result is an instrument you can wear and play.

\section*{The colored boxes: a visual language}

Throughout the book you will meet recurring colored boxes. Each
one marks a specific \emph{kind} of thinking, and learning to read them is part of learning the
book. Here is the legend, and you have already seen several of them in action above.

\begin{whybox}
A \textbf{Why it works} box justifies a technique. When I claim that some procedure -- say,
multiplying two spectra to convolve two signals -- actually does what I say it does, this is
where I show you the reason, usually geometric. If you only read one kind of box, read these:
they explain why the methods work.
\end{whybox}

\begin{happeningbox}
A \textbf{What is really happening} box strips away the notation and tells you, in plain
words, the physical or geometric event behind the math. When an algebraic result looks
surprising, these boxes show the simple mechanism underneath.
\end{happeningbox}

\begin{buildbox}
A \textbf{Build it (hardware)} box is for hands. Wiring notes, sample-rate choices, an
Arduino or Raspberry Pi code sketch, a gotcha about a particular sensor. This is where the
abstractions meet real hardware. You can build the whole device from the contents of these boxes.
\end{buildbox}

\begin{connectionbox}
A \textbf{Connection} box links to the companion volumes -- the ground-up books on
\emph{calculus}, \emph{linear algebra}, and \emph{probability}. Signal processing is where
those three subjects meet: integrals become energy and averaging, inner products and
orthogonality become frequency analysis, and randomness becomes noise. When an idea has a
deeper home in one of those books, I will point you to it, so the four volumes reinforce
rather than repeat each other.
\end{connectionbox}

You will also meet \textbf{Intuition}, \textbf{Formal}, \textbf{History}, \textbf{Common
trap}, \textbf{Application}, \textbf{Try it yourself}, \textbf{Walking question},
\textbf{Why this matters}, and \textbf{Toward the wrist tracker} boxes. Their names say
exactly what they hold. The \emph{History} boxes in particular are worth your time:
knowing that Fourier was nearly laughed out of the French Academy for claiming that any
curve could be built from sines, or that the FFT was rediscovered because the Cold War
needed to detect distant nuclear tests, makes the ideas stick in a way that no derivation
can.

\begin{intuitionbox}
Read the boxes in order the first time through, but on a second pass you can read a chapter
\emph{by box color}. Want only the intuition? Follow the green. Building hardware this
weekend? Skim the orange \textbf{Build it} boxes. Reviewing for an exam? The gray
\textbf{Formal} boxes are your statements of record. The colors are meant as a navigation
system.
\end{intuitionbox}

\section*{How the book is organized}

The book moves in waves, each wave building on the last.

\begin{itemize}
  \item \textbf{Foundations.} What a signal is; continuous versus discrete time; and the
  central reframing that a signal is a \emph{vector} in a space we can do geometry in. This
  is the foundation for everything that follows; do not skip it even if it feels easy.
  \item \textbf{The geometry of frequency.} Sinusoids, phasors, and complex exponentials;
  what frequency really means; the Fourier series, the Fourier transform, the discrete
  Fourier transform, and the FFT; and the short-time Fourier transform that lets us watch
  frequencies change over time. These chapters are the core of the book.
  \item \textbf{Systems and filtering.} Linear time-invariant systems, convolution, impulse
  and frequency response, the $z$- and Laplace transforms, and how to design and apply real
  filters. These are the techniques you will reach for constantly in the build.
  \item \textbf{The bridge to the real world.} Sampling and the Nyquist limit, quantization
  and converters, noise and random signals, spectral estimation, resampling, modulation, and
  a few advanced lenses (wavelets, the analytic signal). This is where we deal with
  the imperfections of real hardware.
  \item \textbf{Domains and the build.} Short tours of audio, images, and motion sensing;
  then a focused march through IMUs, sensing the hand, gesture detection by template
  matching and per-finger calibration, embedded DSP, and finally the \textbf{capstone}:
  assembling the wrist gesture instrument controller end to end.
\end{itemize}

The chapters are written to be read in order, but they are not a chain that snaps if you
jump. If you arrive already knowing the Fourier transform and want filters, go straight to
the systems chapters; the cross-references will send you back for anything you are missing.

\section*{Plots, code, and doing}

This is a visual book and a hands-on book. Expect several figures in every chapter --
waveforms, sums of sinusoids, magnitude and phase spectra, the rotating phasor on the unit
circle, the sinc function, spectrograms, pole-zero plots, filter responses, the sampling
and aliasing picture, convolution sliding past itself, two-dimensional image spectra, and
the jagged traces of real IMU and EMG data. Where a short piece of Python or C makes an idea
concrete, you will find it; signal processing is one of those subjects where ten lines of
code can teach more than a page of prose.

To make the visual language tangible from the very first page, here is a signal and a
glimpse of its frequency content side by side -- the two views of the same object that this
entire book is about.

\begin{figure}[H]
\centering
\begin{tikzpicture}
\begin{axis}[
  name=timeplot,
  width=0.50\textwidth, height=4.6cm,
  title={Time view: a wiggly line},
  title style={font=\small},
  xlabel={\small time $t$}, ylabel={\small amplitude},
  domain=0:1, samples=160,
  xmin=0, xmax=1, ymin=-2.4, ymax=2.4,
  axis lines=box, tick label style={font=\tiny},
  enlargelimits=false,
]
\addplot[sigblue, thick] {sin(deg(2*pi*3*x)) + 0.6*sin(deg(2*pi*7*x))};
\end{axis}
\begin{axis}[
  at={($(timeplot.east)+(1.1cm,0)$)}, anchor=west,
  width=0.44\textwidth, height=4.6cm,
  title={Frequency view: two clean spikes},
  title style={font=\small},
  xlabel={\small frequency (Hz)}, ylabel={\small strength},
  xmin=0, xmax=10, ymin=0, ymax=1.3,
  ytick={0,0.5,1}, xtick={0,3,7,10},
  axis lines=box, tick label style={font=\tiny},
  enlargelimits=false,
]
\addplot[specred, thick, ycomb, mark=*, mark size=1.3pt] coordinates {(3,1.0) (7,0.6)};
\end{axis}
\end{tikzpicture}
\caption{Time and frequency views of one signal. On the left, a signal in the time domain: a single
tangled curve. On the right, the \emph{same} signal seen as the frequencies it is made of --
here, a strong tone at $3$~Hz and a weaker one at $7$~Hz. Learning to move between these two
pictures, in both directions, is the central skill of signal processing.}
\label{ch00:two-views}
\end{figure}

Figure~\ref{ch00:two-views} summarizes the approach of the book. The left curve looks complicated; the
right description is just two numbers and two heights. The Fourier transform is the machine
that carries you from left to right, and the inverse transform carries you back. Most of the
power in this field comes from doing your work in whichever of the two pictures makes it easy,
then traveling to the other.

\begin{labbox}
You do not need any special hardware to start. If you have a computer with Python and NumPy,
you can reproduce nearly every figure in this book and experiment with it. The \textbf{Try it yourself} boxes will suggest small experiments -- generate a
signal, transform it, change one number, watch what moves. You will understand this subject best by
running and modifying code, not only by reading about it.
\end{labbox}

To make the promise of Figure~\ref{ch00:two-views} concrete, look at how the two simple tones
on the right actually \emph{build} the tangled curve on the left. Each tone is a smooth, pure
ripple; the signal is simply their sum. Figure~\ref{ch00:buildup} shows the two
ingredients separately and then stacked. Nothing here requires a formula -- it is addition,
done one column at a time -- yet it already contains the main idea of the Fourier \emph{series}:
complicated things are sums of simple ripples.

\begin{figure}[H]
\centering
\begin{tikzpicture}
\begin{axis}[
  width=0.92\textwidth, height=5.2cm,
  title={Two pure tones, and what you get when you add them},
  title style={font=\small},
  xlabel={\small time $t$ (seconds)}, ylabel={\small amplitude},
  domain=0:1, samples=200,
  xmin=0, xmax=1, ymin=-2.6, ymax=2.6,
  axis lines=box, tick label style={font=\tiny},
  enlargelimits=false,
  legend style={font=\tiny, at={(0.5,-0.28)}, anchor=north, legend columns=3},
]
\addplot[sigblue, thick, dashed] {sin(deg(2*pi*3*x))};
\addplot[specred, thick, dotted] {0.6*sin(deg(2*pi*7*x))};
\addplot[structurepurple, very thick] {sin(deg(2*pi*3*x)) + 0.6*sin(deg(2*pi*7*x))};
\legend{$3$~Hz tone, $7$~Hz tone (weaker), their sum}
\end{axis}
\end{tikzpicture}
\caption{The same signal as in Figure~\ref{ch00:two-views}, taken apart. The dashed curve is a
clean $3$~Hz sinusoid; the dotted curve is a weaker $7$~Hz sinusoid; the solid purple curve is
simply the two added together, point by point. The ``wiggly line'' is a
recipe: two tones in fixed proportions. Reading that recipe off the curve is what the Fourier transform does for us.}
\label{ch00:buildup}
\end{figure}

\begin{happeningbox}
Look again at the solid curve in Figure~\ref{ch00:buildup} and try to \emph{see} both tones at
once: the slow, three-humps-per-second swell of the $3$~Hz tone, with the faster $7$~Hz ripple
riding on top of it like texture on a wave. Once your eye learns to do this -- to perceive a
single curve as a stack of pure tones -- you have begun to think in the frequency domain
without any mathematics at all. The transform is just the instrument that makes this perception
exact and reversible.
\end{happeningbox}

A second habit this book will build is reading a \emph{spectrum} -- the right-hand picture --
as fluently as you read a waveform. A spectrum answers one question: \emph{how much of each
frequency is present?} Tall spike, lots of that frequency; no spike, none of it. Here is a
worked reading of one, the kind we will do constantly.

\begin{example}[Reading a magnitude spectrum]
Suppose an accelerometer on the wrist, sampled while the hand performs a steady tremor plus a
slow deliberate wave, produces the magnitude spectrum sketched in Figure~\ref{ch00:readspec}.
What can we say \emph{just by looking}?

\begin{enumerate}
  \item There is a tall spike near $1$~Hz. A once-per-second oscillation is far too slow to be
  tremor; this is the \emph{deliberate} hand wave. Its height tells us it dominates the motion's
  energy.
  \item There is a second, shorter cluster near $6$~Hz. Physiological hand tremor famously lives
  in roughly the $4$ to $12$~Hz band, so this bump is almost certainly involuntary tremor.
  \item Everything above $10$~Hz is essentially flat and small -- the noise floor of the sensor,
  with no real signal sticking out of it.
\end{enumerate}

The design consequence follows directly: to recover the intentional gesture and suppress the
tremor, we want a filter that \emph{keeps} the low-frequency spike and \emph{rejects} the
$6$~Hz bump. We have not derived a single filter equation yet, and we already know what the
filter must do, because we read it off the spectrum. That is the practical payoff of thinking in frequency.
\end{example}

\begin{figure}[H]
\centering
\begin{tikzpicture}
\begin{axis}[
  width=0.92\textwidth, height=5.0cm,
  title={A magnitude spectrum you can read at a glance},
  title style={font=\small},
  xlabel={\small frequency (Hz)}, ylabel={\small strength},
  domain=0:20, samples=200,
  xmin=0, xmax=20, ymin=0, ymax=1.25,
  xtick={0,1,6,10,15,20}, ytick={0,0.5,1},
  axis lines=box, tick label style={font=\tiny},
  enlargelimits=false,
]
% deliberate wave: tall narrow peak near 1 Hz
\addplot[name path=curve, structurepurple, very thick, smooth]
  {1.0*exp(-((x-1)^2)/0.18) + 0.42*exp(-((x-6)^2)/1.4) + 0.06};
\path[name path=axis] (axis cs:0,0) -- (axis cs:20,0);
\addplot[fillsig, opacity=0.35] fill between[of=curve and axis];
\node[font=\tiny, sigblue, anchor=south] at (axis cs:1.0,1.10) {deliberate wave};
\node[font=\tiny, specred, anchor=south] at (axis cs:6.0,0.52) {tremor band};
\node[font=\tiny, gray, anchor=south west] at (axis cs:14,0.10) {noise floor};
\end{axis}
\end{tikzpicture}
\caption{A made-up but realistic magnitude spectrum of wrist motion. A tall, narrow peak at low
frequency is the intentional gesture; a broader, shorter bump in the $4$--$12$~Hz region is
involuntary tremor; the flat tail is sensor noise. Reading these three features -- and deciding
what to keep and what to throw away -- is the everyday work of the signal processor.}
\label{ch00:readspec}
\end{figure}

\section*{A note on rigor}

I have tried to make this book both intuitive and correct, and where those two goals pull
apart I handle the tension openly rather than quietly fudging. When I simplify, I will tell
you I am simplifying and where the full story lives. When a statement is true only under
conditions, the conditions appear in a \textbf{Formal} box rather than being swept away. The
geometry comes first because it is how human beings actually understand; the rigor comes
second because it is what keeps the intuition from lying to us.

The same discipline governs the engineering. Real sensors are noisy, drift with temperature,
and shift on the body. Real microcontrollers run out of cycles and memory. Real wireless
links drop packets at the worst moment. I will not pretend these problems away to make the
build sound easier than it is. A device that works in the messy real world is worth more than
a clean idea that only works in a slide deck, and getting there means dealing with those problems directly.

\section*{Acknowledgment of the giants}

\begin{historybox}
Signal processing grew out of concrete engineering problems and mathematical
debates. From the 1850s through the 1880s, laying the first
submarine telegraph cables forced \textbf{William Thomson (Lord Kelvin)} and \textbf{Oliver
Heaviside} to model how undersea conductors smear sharp pulses and how distributed inductance
and operational calculus can tame that distortion -- a story we tell in Chapter~1. Over the
next century, the same pressure to push clean signals through noisy physical channels drove
breakthroughs in long-distance telephony, radar, seismic monitoring, and digital audio. Each
time hardware hit a physical limit, new mathematics had to be invented or sharpened to get past
it.

The key mathematical ideas came from many people over two centuries:
\textbf{Leonhard Euler} (1748) and \textbf{Daniel Bernoulli} debated the harmonics of vibrating
strings; \textbf{Joseph Fourier} (1807) braved the fierce skepticism of Lagrange and Laplace to
argue that any wave is a sum of sinusoids; \textbf{Carl Friedrich Gauss} (1805) quietly invented
the fast Fourier transform to track the asteroid Pallas 160 years before \textbf{James Cooley}
and \textbf{John Tukey} (1965) rediscovered it to detect Soviet underground nuclear detonations;
\textbf{Harry Nyquist} (1928) and \textbf{Claude Shannon} (1948) proved the fundamental limits of
sampling and information; and \textbf{Dennis Gabor} (1946) established the quantum-analogous
time-frequency uncertainty principle. This book builds on their work.
\end{historybox}

\begin{connectionbox}
\textbf{How the Book Builds: Each Tool Solves One Problem and Exposes the Next}

Each part of this book answers a question left open by the part before it:
\begin{itemize}
  \item \textbf{Part I (Chapters 1--4) -- Signals, Vectors, and Phasors}: Treats signals as vectors
  so we can measure energy, similarity, and orthogonal projections geometrically, then introduces
  rotating complex arrows (phasors) to turn trig identities and derivatives into simple algebra.
  \emph{What is still missing}: Looking at a wiggly time trace still does not tell you which
  frequencies are mixed inside it.
  \item \textbf{Part II (Chapters 5--11) -- Frequency and the Fourier Family}: Projects signals onto
  pure sinusoids to read off their frequency recipes -- via Fourier series, the continuous Fourier
  transform, the discrete Fourier transform (DFT), and the fast $\mathcal{O}(N\log N)$ FFT.
  \emph{What is still missing}: A single Fourier transform averages over the whole recording, telling
  you \emph{which} notes were played but losing \emph{when} each note happened. Sliding a short window
  gives the spectrogram (Chapter~10), at the price of a fixed tradeoff between time sharpness and
  frequency sharpness.
  \item \textbf{Part III (Chapters 12--17) -- Systems and Filtering}: Shows why linear time-invariant
  (LTI) systems pass every sinusoid through at the same frequency, letting us sculpt spectra using
  convolution and pole--zero geometry in the $z$-plane. \emph{What is still missing}: Ideal brick-wall
  filters would have to see into the future and ring forever, so real hardware filters must balance
  sharpness, delay, and phase distortion.
  \item \textbf{Parts IV \& V (Chapters 18--25) -- Sampling, Noise, and Multirate}: Connects continuous
  physics to finite digital samples via the Nyquist--Shannon sampling theorem, models quantization and
  random noise, and introduces wavelets (Chapter~24), whose window stretches and shrinks with frequency.
  \item \textbf{Parts VI--VIII (Chapters 26--41) -- Audio, Images, and Wearable Embedded DSP}: Puts
  every tool to work on sound, 2D images, IMU orientation fusion, forearm muscle sensing, and the
  real-time C firmware of our wrist gesture controller.
\end{itemize}
\end{connectionbox}

\section*{Warm-up explorations}

A preface is an odd place for exercises, so think of these less as homework and more as a
warm-up. None requires anything you have not already read
on these pages; all of them reward five minutes at a keyboard or with a pencil. If they feel
easy, good -- they are meant to. If they feel hard, they point to the skills this book
will build.

\begin{exercise}[Add two tones by hand]
On graph paper (or in your head), sketch one full second of a $2$~Hz sine wave and, beneath it,
a $5$~Hz sine wave at half the height. Now add them sample by sample at a few points -- say every
tenth of a second -- and plot the sum. Compare your sketch to the spirit of
Figure~\ref{ch00:buildup}. Which tone sets the overall ``shape,'' and which adds the ``texture''?
\end{exercise}

\begin{exercise}[Read a spectrum cold]
Cover the labels in Figure~\ref{ch00:readspec} and describe, in one sentence each, what the
three features are and what a designer would want to do with them. Then uncover the labels and
check yourself. The goal is to make spectrum-reading as automatic as reading a thermometer.
\end{exercise}

\begin{exercise}[Name the box]
Flip back through this preface and, for each colored box you met, write a single phrase naming
the \emph{kind of thinking} it represents (justification, plain-language picture, hardware,
limits, and so on). This helps you learn the book's navigation system.
\end{exercise}

\begin{exercise}[State the sensor contract]
In your own words, and without looking, explain (a) why \emph{which} finger moved \emph{is}
detectable from a forearm FMG or EMG band given per-finger calibration, and (b) why the
\emph{continuous joint angle} of each finger is \emph{not} recoverable from a wrist IMU.
Name a sensor that would close the second gap. If you can explain this to a friend, you understand the
sensor limits that the rest of the build depends on.
\end{exercise}

\begin{exercise}[Your first FFT]
If you have Python and NumPy, generate the very signal of Figure~\ref{ch00:two-views} --
\texttt{np.sin(2*np.pi*3*t) + 0.6*np.sin(2*np.pi*7*t)} on a fine time grid \texttt{t} -- and
plot it. Then call \texttt{np.fft.rfft} on it and plot the magnitude. You will not understand
the FFT yet, and that is fine; the point is to \emph{see}, with your own data, the two spikes
appear where the book promised they would. Keep this script. We will return to it.
\end{exercise}

\section*{How to begin}

Turn the page. Read slowly at first; the early chapters introduce ideas that are used much later, and
time spent making the foundations \emph{feel} obvious is repaid many times over. Keep a Python
session open. When a box says \emph{try it}, try it. When the geometry clicks -- when you first
\emph{see} the rotating phasor, or watch a filter erase a hum from a recording, or read a tap
straight off a wiggling accelerometer trace -- you will see why this subject is worth the effort.

Welcome to the geometry of frequency. Chapter~1 begins with the most basic question: what is a signal?
